\documentclass[english]{article}

\usepackage{ae,aecompl}
\usepackage[T1]{fontenc}
\usepackage[latin9]{inputenc}
\usepackage{textcomp}
\usepackage{graphicx}
\usepackage{eurosym}
\usepackage{hyperref}
\usepackage{float}
\usepackage{tabto}

\makeatletter
\usepackage{babel}


\makeatother

\usepackage{babel}
\begin{document}
\begin{figure}
	\centering
	\includegraphics[scale=0.5]{logo.pdf} 
\end{figure}


\title{PowerEnJoy\\
 Requirement Analysis and Specification Document\\
}

\date{A.A 2016/2017}

\author{Erba Alessandro\\
 Leveni Filippo\\
 Lodi Luca}

\maketitle
\pagebreak{}

\tableofcontents{} \pagebreak{}

\section{Introduction}
	\subsection{Purpose }
	
		The purpose of our team is to project PowerEnjoy, which is  a platform that will be used to manage a car sharing service 
		in the city area.\\
		The platform will allow clients to find and reserve an electric car using mobile app from their self phone, and it will also 			manage the communication of the emergency situations to the proper staff.\\
		Furthermore, a particular PowerEnjoy feature is the application of certain discounts to its users.\\
		\tab \tab This document aims to describe the high-level functionalities that will be offered by PowerEnjoy service. The 			RASD is intended to be viewed by the stakeholders, that will evaluate the correctness of any assumption and decision
		in this document.
		
	\subsection{Scope and Problem Description}
	
		The government of a large city would like to introduce an ecological way to travel within the city area. In particular, 				uniform distribution of cars in the city is a significant target.\\
		\tab \tab The cars are electric as anticipated and inside them it is available an integrated interface. Through this interface			is possible to communicate with the system (e.g. to call emergency staff).\\
		\tab \tab Users can choose and reserve a car through a mobile app. The system answers to the request by showing the 			time left to the expiration of the reservation.\\
		The staff is divided in three categories:\\
			- Field staff\\
			- Emergency staff\\
			- Management staff\\
		\tab \tab Field staff use the mobile application to inform the system about their availability and to confirm that they are 			going to take care of a certain call (e.g. retrieve a broken car, move a vehicle with a low battery or to ensure the uniform 			distribution in the city, etc...).\\
		\tab \tab Emergency staff use the mobile application to answer customers calls and to take care of special situations (e.g. 			no money in the costumer's credit card, car accident, etc...).\\
		\tab \tab Management staff use the mobile application to handle the service (e.g. set the safe areas, fares, discounts, 				etc...).\\
		\tab \tab The system guarantees a uniform distribution of cars in the city. In particular, the city is divided in 						circumferences and sectors. The system automatically computes the distribution of available cars both in the areas 				between the circumferences and in the sectors, based on the GPS information it receives from each car. Uniformity is 				guaranteed by ensuring the same amount of available cars for each area between the circumferences and for each sector.			\\
		\tab \tab A user can take advantage of the "Money saving" option selecting it on the integrated interface in the car. After 			that he/she can input his/her final destination and the system provides information about the station where to leave the 			car to get a discount. This destination is determined to ensure the uniform distribution of cars and depends both 				on the destination of the user and on the availability of power plugs at the selected station.
	
	\subsection{Goals}
	\subsection{Domain Assumptions }
	\subsection{Glossary}
	\subsection{Further Developments }
	\subsection{Used Tools}

\section{Specific Requirements}
	\subsection{Functional Requirements}
	\subsection{Non Functional Requirements}

\section{Scenarios Identifying}
\section{UML Models And Use Cases}
	\subsection{Use Cases Diagram}
	\subsection{Actors Identifying}
	\subsection{Use Cases}

	\subsection{Class Diagram}
\section{External Interfaces}
\section{Alloy Model}
\section{Hours of Work}

\end{document}